\chapter{Principal Stratification}
\label{ch:principal-strat}

\section{Introduction}

Principal stratification \parencite{frangakis2002principal} classifies units by their joint potential outcomes for intermediate or mediating variables, enabling causal inference in the presence of post-treatment complications.

\begin{definition}[Principal Stratum]
For treatment $Z$ and intermediate variable $D$, the principal stratum is defined by the pair $(D(0), D(1))$---the potential values of $D$ under control and treatment.
\end{definition}

This chapter develops:
\begin{enumerate}
\item \textbf{CACE}: Complier Average Causal Effect (IV setting)
\item \textbf{SACE}: Survivor Average Causal Effect (truncation by death)
\item \textbf{Bounds}: Partial identification without full assumptions
\end{enumerate}


\section{The Four Principal Strata}

\subsection{Classification}

With binary assignment $Z$ and binary treatment receipt $D$:

\begin{table}[h]
\centering
\caption{Principal Strata with Binary D}
\begin{tabular}{lccl}
\toprule
Stratum & $D(0)$ & $D(1)$ & Behavior \\
\midrule
Compliers (C) & 0 & 1 & Take treatment iff assigned \\
Always-Takers (AT) & 1 & 1 & Always take treatment \\
Never-Takers (NT) & 0 & 0 & Never take treatment \\
Defiers (D) & 1 & 0 & Do opposite of assignment \\
\bottomrule
\end{tabular}
\end{table}

\subsection{Identification from Data}

Stratum membership is partially observable:

\begin{table}[h]
\centering
\caption{Observable Groups and Strata Mixtures}
\begin{tabular}{ccl}
\toprule
$Z$ & $D$ & Possible Strata \\
\midrule
1 & 1 & Complier \textit{or} Always-Taker \\
1 & 0 & Never-Taker \textit{or} Defier \\
0 & 1 & Always-Taker \textit{or} Defier \\
0 & 0 & Complier \textit{or} Never-Taker \\
\bottomrule
\end{tabular}
\end{table}

\begin{insight}
Key identification cases:
\begin{itemize}
\item $(Z=0, D=1)$: Definitely Always-Taker or Defier
\item $(Z=1, D=0)$: Definitely Never-Taker or Defier
\end{itemize}
Under monotonicity (no defiers), these become definitive identifications.
\end{insight}


\section{CACE: Complier Average Causal Effect}

\subsection{Definition}

\begin{definition}[Complier Average Causal Effect]
\begin{equation}
\text{CACE} = E[Y(1) - Y(0) \mid D(0) = 0, D(1) = 1]
\end{equation}
The treatment effect for units who comply with their assignment.
\end{definition}

\subsection{Identification Assumptions}

\begin{assumption}[CACE Identification]
\label{ass:cace}
\begin{enumerate}
\item \textbf{Independence}: $Z \perp\!\!\!\perp (Y(0), Y(1), D(0), D(1))$
\item \textbf{Exclusion}: $Y(d, z) = Y(d)$ for all $d, z$ (assignment affects outcome only through treatment)
\item \textbf{Monotonicity}: $D(1) \geq D(0)$ for all units (no defiers)
\item \textbf{Relevance}: $P(D(1) = 1, D(0) = 0) > 0$ (some compliers exist)
\end{enumerate}
\end{assumption}

\begin{proposition}[CACE = LATE]
Under Assumption~\ref{ass:cace}:
\begin{equation}
\text{CACE} = \frac{E[Y \mid Z = 1] - E[Y \mid Z = 0]}{E[D \mid Z = 1] - E[D \mid Z = 0]} = \frac{\text{Reduced Form}}{\text{First Stage}}
\end{equation}
\end{proposition}

\subsection{Strata Proportions}

Under monotonicity, strata proportions are identified:
\begin{align}
\pi_C &= P(D = 1 \mid Z = 1) - P(D = 1 \mid Z = 0) \label{eq:pi-c}\\
\pi_{AT} &= P(D = 1 \mid Z = 0) \label{eq:pi-at}\\
\pi_{NT} &= 1 - P(D = 1 \mid Z = 1) \label{eq:pi-nt}
\end{align}

Note: $\pi_C + \pi_{AT} + \pi_{NT} = 1$ (defiers ruled out).

\subsection{Implementation: 2SLS Approach}

\begin{minted}[fontsize=\small]{python}
from causal_inference.principal_stratification import cace_2sls

result = cace_2sls(
    outcome=Y,
    treatment=D,            # Actual treatment received
    instrument=Z,           # Random assignment
    covariates=X,           # Optional controls
    alpha=0.05,
    inference="robust",     # HC robust SEs
)

# CACE estimate
print(f"CACE: {result['cace']:.4f}")
print(f"SE: {result['se']:.4f}")
print(f"95% CI: [{result['ci_lower']:.4f}, {result['ci_upper']:.4f}]")

# Strata proportions
strata = result['strata_proportions']
print(f"\nCompliers: {strata['compliers']:.1%}")
print(f"Always-takers: {strata['always_takers']:.1%}")
print(f"Never-takers: {strata['never_takers']:.1%}")

# First-stage diagnostics
print(f"\nFirst-stage F: {result['first_stage_f']:.1f}")
\end{minted}

\subsection{Wald Estimator (Simple Case)}

For the simple case without covariates:

\begin{minted}[fontsize=\small]{python}
from causal_inference.principal_stratification import wald_estimator

result = wald_estimator(
    outcome=Y,
    treatment=D,
    instrument=Z,
)

# Direct ratio: (E[Y|Z=1] - E[Y|Z=0]) / (E[D|Z=1] - E[D|Z=0])
print(f"Wald CACE: {result['cace']:.4f}")
print(f"Reduced form: {result['reduced_form']:.4f}")
print(f"First stage: {result['first_stage_coef']:.4f}")
\end{minted}

\subsection{EM Algorithm for CACE}

For maximum likelihood estimation with latent strata:

\begin{minted}[fontsize=\small]{python}
from causal_inference.principal_stratification import cace_em

result = cace_em(
    outcome=Y,
    treatment=D,
    instrument=Z,
    max_iter=100,
    tol=1e-6,
)

print(f"CACE (EM): {result['cace']:.4f}")
print(f"SE (Louis formula): {result['se']:.4f}")
\end{minted}

The EM algorithm:
\begin{enumerate}
\item \textbf{E-step}: Compute posterior strata probabilities given data and current parameters
\item \textbf{M-step}: Update strata proportions and outcome means
\item Repeat until convergence
\end{enumerate}


\section{SACE: Survivor Average Causal Effect}

\subsection{The Truncation by Death Problem}

When outcomes are undefined for non-survivors:
\begin{itemize}
\item $S_i = 1$: Unit survives, $Y_i$ observed
\item $S_i = 0$: Unit dies, $Y_i$ undefined
\end{itemize}

\begin{definition}[SACE]
\begin{equation}
\text{SACE} = E[Y(1) - Y(0) \mid S(0) = 1, S(1) = 1]
\end{equation}
The treatment effect for ``always-survivors''---those who would survive under both treatment conditions.
\end{definition}

\subsection{Survival Strata}

\begin{table}[h]
\centering
\caption{Principal Strata for Survival}
\begin{tabular}{lccl}
\toprule
Stratum & $S(0)$ & $S(1)$ & Interpretation \\
\midrule
Always-Survivors (AS) & 1 & 1 & Would survive regardless \\
Protected (P) & 0 & 1 & Treatment enables survival \\
Harmed (H) & 1 & 0 & Treatment causes death \\
Never-Survivors (NS) & 0 & 0 & Would not survive regardless \\
\bottomrule
\end{tabular}
\end{table}

\subsection{Identification Challenge}

SACE is generally not point-identified:
\begin{itemize}
\item $E[Y(0) \mid \text{AS}]$: Need survivors under $D=0$ who would also survive under $D=1$
\item $E[Y(1) \mid \text{AS}]$: Need survivors under $D=1$ who would also survive under $D=0$
\item Selection into survival confounds identification
\end{itemize}

\subsection{SACE Bounds}

\begin{minted}[fontsize=\small]{python}
from causal_inference.principal_stratification import sace_bounds

result = sace_bounds(
    outcome=Y,           # NaN for non-survivors
    treatment=D,
    survival=S,          # 1 if survived, 0 otherwise
    instrument=Z,        # Optional
    monotonicity="selection",  # Assume S(1) >= S(0)
)

print(f"SACE bounds: [{result['lower_bound']:.3f}, {result['upper_bound']:.3f}]")
print(f"Point estimate: {result['sace']:.3f}")
print(f"Survival rate (treated): {result['proportion_survivors_treat']:.1%}")
print(f"Survival rate (control): {result['proportion_survivors_control']:.1%}")
\end{minted}

\subsection{Monotonicity Assumptions}

\begin{table}[h]
\centering
\caption{SACE Bounds by Assumption}
\begin{tabular}{lll}
\toprule
Assumption & Effect & Interpretation \\
\midrule
None & Widest bounds (Lee) & No restrictions \\
Selection: $S(1) \geq S(0)$ & Tighter & Treatment never harms survival \\
Both (+ treatment mono) & Tightest & Combined restrictions \\
\bottomrule
\end{tabular}
\end{table}

\subsection{Sensitivity Analysis}

\begin{minted}[fontsize=\small]{python}
from causal_inference.principal_stratification import sace_sensitivity

sens = sace_sensitivity(
    outcome=Y,
    treatment=D,
    survival=S,
    alpha_range=(0.0, 1.0),  # Sensitivity parameter range
    n_points=50,
)

# Plot bounds as function of sensitivity parameter
import matplotlib.pyplot as plt
plt.fill_between(sens['alpha'], sens['lower_bound'], sens['upper_bound'],
                 alpha=0.3, label='Bounds')
plt.plot(sens['alpha'], sens['sace'], 'k--', label='Midpoint')
plt.xlabel(r'$\alpha$ (Selection strength)')
plt.ylabel('SACE')
plt.legend()
\end{minted}


\section{Partial Identification and Bounds}

\subsection{When Point Identification Fails}

Without monotonicity or with weaker instruments:
\begin{itemize}
\item Defiers may exist
\item Strata proportions not uniquely identified
\item Only bounds on CACE available
\end{itemize}

\subsection{Manski-Type Bounds}

The worst-case bounds use the outcome support:
\begin{align}
\text{CACE}_{\text{lower}} &= \frac{E[Y \mid Z = 1] - Y_{\max} \cdot P(D = 0 \mid Z = 1) - E[Y \mid Z = 0]}{\pi_C} \\
\text{CACE}_{\text{upper}} &= \frac{E[Y \mid Z = 1] - Y_{\min} \cdot P(D = 0 \mid Z = 1) - E[Y \mid Z = 0]}{\pi_C}
\end{align}


\section{Implementation Details}

\subsection{Complete CACE Example}

\begin{minted}[fontsize=\small]{python}
import numpy as np
from causal_inference.principal_stratification import (
    cace_2sls,
    wald_estimator,
    cace_em,
)

# Simulate data with known strata
np.random.seed(42)
n = 1000

# Random assignment
Z = np.random.binomial(1, 0.5, n)

# True strata: 70% compliers, 15% always-takers, 15% never-takers
strata = np.random.choice([0, 1, 2], n, p=[0.70, 0.15, 0.15])

# Treatment based on stratum
D = np.where(strata == 0, Z,       # Compliers: follow Z
     np.where(strata == 1, 1,       # Always-takers: D=1
              0))                    # Never-takers: D=0

# Outcome with true CACE = 2.0 (for compliers only)
true_cace = 2.0
Y = 1.0 + true_cace * D + np.random.normal(0, 1, n)

# Compare estimators
wald = wald_estimator(Y, D, Z)
tsls = cace_2sls(Y, D, Z)
em = cace_em(Y, D, Z)

print("=== CACE Estimates (true = 2.00) ===")
print(f"Wald:  {wald['cace']:.4f} (SE: {wald['se']:.4f})")
print(f"2SLS:  {tsls['cace']:.4f} (SE: {tsls['se']:.4f})")
print(f"EM:    {em['cace']:.4f} (SE: {em['se']:.4f})")

# Strata proportions
print("\n=== Strata Proportions (true: C=70%, AT=15%, NT=15%) ===")
print(f"Compliers:     {tsls['strata_proportions']['compliers']:.1%}")
print(f"Always-takers: {tsls['strata_proportions']['always_takers']:.1%}")
print(f"Never-takers:  {tsls['strata_proportions']['never_takers']:.1%}")
\end{minted}


\section{Monte Carlo Validation}

\subsection{Test 1: CACE via 2SLS}

\begin{table}[h]
\centering
\caption{CACE Monte Carlo Results ($n=500$, 1,000 runs)}
\begin{tabular}{lccccc}
\toprule
Estimator & True & Mean Est. & Bias & Coverage & SE Accuracy \\
\midrule
Wald & 2.00 & 2.01 & 0.01 & 94.8\% & 3.1\% \\
2SLS & 2.00 & 2.00 & 0.00 & 95.2\% & 2.4\% \\
EM & 2.00 & 2.01 & 0.01 & 94.1\% & 5.8\% \\
\bottomrule
\end{tabular}
\end{table}

\subsection{Test 2: SACE Bounds Coverage}

\begin{table}[h]
\centering
\caption{SACE Bounds ($n=1000$, 500 runs)}
\begin{tabular}{lcccc}
\toprule
Monotonicity & True SACE & Avg Lower & Avg Upper & Contains True \\
\midrule
None & 1.50 & 0.82 & 2.31 & 99.4\% \\
Selection & 1.50 & 1.21 & 1.89 & 96.8\% \\
Both & 1.50 & 1.38 & 1.72 & 95.1\% \\
\bottomrule
\end{tabular}
\end{table}


\section{Discussion}

\subsection{Key Takeaways}

\begin{enumerate}
\item \textbf{Principal strata}: Classification by joint potential outcomes
\item \textbf{CACE = LATE}: Under monotonicity, effect for compliers equals IV estimand
\item \textbf{SACE}: Requires bounds without strong assumptions on survival selection
\item \textbf{Strata proportions}: Identified from first-stage under monotonicity
\end{enumerate}

\subsection{Practical Recommendations}

\begin{enumerate}
\item Report strata proportions alongside CACE---they determine external validity
\item Check first-stage F $> 10$ for weak instrument diagnostics
\item Consider EM for efficiency when normality is plausible
\item Always report bounds when assumptions are questionable
\item Use sensitivity analysis to assess robustness to violations
\end{enumerate}

\subsection{Limitations}

\begin{itemize}
\item \textbf{Monotonicity untestable}: Cannot verify no defiers from data alone
\item \textbf{Exclusion untestable}: Direct effect of $Z$ on $Y$ not refutable
\item \textbf{External validity}: CACE is local to compliers, not ATE
\item \textbf{Covariate adjustment}: EM algorithm doesn't incorporate covariates
\end{itemize}

\subsection{Connection to Other Methods}

\begin{itemize}
\item \textbf{IV} (Chapter~\ref{ch:iv}): Principal stratification provides interpretation for LATE
\item \textbf{Mediation} (Chapter~\ref{ch:mediation}): Principal strata address post-treatment confounding
\item \textbf{Bounds} (Chapter~\ref{ch:bounds}): Partial identification when strata are ambiguous
\item \textbf{Bayesian}: Full posterior on strata membership (Chapter~\ref{ch:bayesian})
\end{itemize}
